\section{АНАЛИЗ ПРЕДМЕТНОЙ ОБЛАСТИ}

\subsection{Обзор существующих программных продуктов по теме работы}

На современном рынке программного обеспечения для управления проектами представлено множество решений. Для выявления недостатков существующих систем и формирования требований к разрабатываемому веб-приложению проведен структурированный анализ популярных платформ: Jira, Trello и Asana. 

Сравнение проводилось по ключевым критериям, критически важным для небольших рабочих групп: наличие полнофункциональной бесплатной версии, простота освоения, поддержка ролевой модели доступа (RBAC), наличие встроенной статистики и возможность прикрепления файлов. Результаты сравнительного анализа представлены в таблице \ref{tab:competitors}.

\begin{table}[h!]
\centering
\caption{Сравнительный анализ систем управления проектами}
\label{tab:competitors}
\begin{tabular}{|p{0.3\textwidth}|c|c|c|c|}
\hline
\textbf{Критерий оценки} & \textbf{Jira} & \textbf{Trello} & \textbf{Asana} & \textbf{Наш проект} \\
\hline
Бесплатный полный функционал & -- & -- & -- & + \\
\hline
Простота освоения & -- & + & + & + \\
\hline
Ролевая модель доступа & + & -- (платно) & -- (платно) & + \\
\hline
Графическая статистика & + & -- & + & + \\
\hline
Прикрепление файлов & + & + & + & + \\
\hline
\end{tabular}
\end{table}

Анализ показал, что мощные корпоративные системы (Jira) обладают необходимым функционалом, но перегружены и сложны для малого бизнеса. Более простые решения (Trello, Asana) искусственно ограничивают функционал в бесплатных версиях, вынуждая покупать подписку ради доступа к статистике или ролевым моделям. Это подтверждает целесообразность разработки собственного легковесного веб-приложения.

\subsection{Анализ программных инструментов разработки веб-приложений}

Для реализации клиент-серверной архитектуры проекта необходимо выбрать оптимальный стек технологий. Сравнение производилось по трем ключевым направлениям: серверная часть (Backend), база данных и клиентская часть (Frontend).

\textbf{1. Выбор Backend-фреймворка.} Рассматривались популярные решения на языке Python: \textit{Django}, \textit{FastAPI} и \textit{Flask}.
\begin{itemize}
    \item \textit{Django} отличается монолитной архитектурой «всё включено», имеет встроенную ORM и панель администратора. Однако его избыточность негативно сказывается на скорости разработки небольших проектов.
    \item \textit{FastAPI} выделяется асинхронной архитектурой и высокой производительностью, идеально подходит для построения микросервисов (API). Недостаток — сложность интеграции с классическим серверным рендерингом страниц.
    \item \textit{Flask} \cite{flask_doc} предоставляет минималистичное ядро с возможностью точечного подключения нужных библиотек (высокая модульность).
\end{itemize}

\textbf{Вывод:} Фреймворк Flask выбран как оптимальный баланс между легковесностью и функциональностью, позволяющий реализовать приложение без избыточного кода.

\textbf{2. Выбор системы управления базами данных (СУБД).} Сравнение проводилось между \textit{SQLite}, \textit{MySQL} и \textit{PostgreSQL}.
\begin{itemize}
    \item \textit{SQLite} хранит данные в одном локальном файле, не требует настройки сервера, но блокирует базу целиком при записи, что исключает её использование в многопользовательской среде.
    \item \textit{MySQL} обладает высокой скоростью чтения, однако исторически имеет менее строгую проверку типов данных и сложную систему изменения таблиц.
    \item \textit{PostgreSQL} является мощной объектно-реляционной СУБД, поддерживающей механизм многоверсионности (MVCC) для параллельной работы без блокировок, строгую ссылочную целостность и сложные аналитические запросы.
\end{itemize}

\textbf{Вывод:} Для надежного хранения данных пользователей и работы с агрегацией статистики в production-среде выбрана СУБД PostgreSQL, взаимодействующая с кодом через ORM-библиотеку SQLAlchemy \cite{sqlalchemy_doc}.

\textbf{3. Выбор Frontend-технологий.} Сравнивались подходы Single Page Application (SPA) и Server-Side Rendering (SSR).
\begin{itemize}
    \item Использование SPA-фреймворков (\textit{React} или \textit{Vue.js}) позволяет создать высокоинтерактивный интерфейс без перезагрузки страниц, но требует настройки отдельного сервера сборки (Node.js/Webpack) и усложняет маршрутизацию запросов (CORS).
    \item Подход SSR на базе шаблонизатора \textit{Jinja2} позволяет генерировать HTML-страницы непосредственно на Python-сервере. В связке с CSS-библиотекой \textit{Bootstrap 5} это обеспечивает быструю отрисовку и адаптивность под мобильные устройства «из коробки».
\end{itemize}

\textbf{Вывод:} Для интерфейса Task Tracker'а динамики связки Jinja2 + Bootstrap 5 более чем достаточно, что исключает необходимость в избыточной SPA-архитектуре.

\subsection{Формулировка цели и задач работы}

На основе проведенного анализа предметной области (п. 1.1) и выбора технологического стека (п. 1.2), главная \textbf{цель работы} заключается в создании веб-приложения для управления проектами и отслеживания задач на базе фреймворка Flask и СУБД PostgreSQL.

Для достижения поставленной цели необходимо решить следующие \textbf{задачи}:
\begin{enumerate}
    \item Провести анализ предметной области и существующих решений;
    \item Спроектировать пользовательские сценарии (Use Cases) и проанализировать целевую аудиторию;
    \item Спроектировать логическую и физическую модели реляционной базы данных;
    \item Разработать макеты пользовательского интерфейса (Wireframes);
    \item Реализовать подсистему аутентификации пользователей и настроить систему распределения ролей;
    \item Разработать программные модули (CRUD-интерфейсы) для управления сущностями «Проект», «Задача» и «Пользователь»;
    \item Реализовать функционал загрузки и скачивания файлов, прикрепленных к задачам;
    \item Разработать модуль агрегирования данных для вывода аналитической статистики по проектам;
    \item Развернуть готовое веб-приложение на удаленном сервере (хостинге).
\end{enumerate}